\documentclass[12pt,a4paper]{leaflet}
\usepackage[utf8]{inputenc}
\usepackage[german]{babel}
\usepackage[T1]{fontenc}
\usepackage{microtype}
\usepackage{graphicx}
\usepackage{color}
\definecolor{gray}{gray}{.5}
\definecolor{darkgray}{gray}{.3}
\setlength{\parindent}{0pt} 
\usepackage{blindtext}
\usepackage[condensed]{tgheros}
\renewcommand*\familydefault{\sfdefault}
\setmargins{2cm}{1cm}{8mm}{8mm}
\usepackage[absolute]{textpos}

\usepackage{hyperref}
\hypersetup{
    pdftoolbar=true,
    pdfmenubar=true,
    pdffitwindow=false,
    pdfstartview={FitH},
    pdfnewwindow=true,
    colorlinks=true,
    linkcolor=black,
    citecolor=black,
    filecolor=black,
    urlcolor=black
}
\def\UrlFont{\bfseries}
\pagestyle{empty} 
\AddToBackground*{2}{\put(0,25){\includegraphics[scale=1]{welle2}}}

%%% Softwaretemplate %%%
\newcommand{\software}[7]{
\begin{minipage}[t][0.4\textheight]{1\textwidth}
\flushleft
\section*{#1 \hfill \normalfont\textcolor{gray}{#2}}
#3 
\medskip

%\scalebox{.9}[1.0]{\textit{Betriebssystem:} \textbf{#4}} \\
\textit{Betriebssystem:} \textbf{#4} \\
\textit{Lizenz:} \textbf{#6} \\
\textit{Deutschsprachig:} \textbf{#7} \\
\textit{Homepage:} \url{#5} \\
\vglue  2 true cm
\end{minipage}
}

% Hack for LibreOffice/OpenOffice: Version of \software without link
\newcommand{\softwaretwo}[7]{
\begin{minipage}[t][0.4\textheight]{1\textwidth}
\flushleft
\section*{#1 \hfill \normalfont\textcolor{gray}{#2}}
#3 
\medskip

%\scalebox{.9}[1.0]{\textit{Betriebssystem:} \textbf{#4}} \\
\textit{Betriebssystem:} \textbf{#4} \\
\textit{Lizenz:} \textbf{#6} \\
\textit{Deutschsprachig:} \textbf{#7} \\
\textit{Homepage:} \textbf{#5} \\
\vglue  2 true cm
\end{minipage}
}


%%% Frontpage %%%
\newcommand{\frontpage}[1]{

\begin{titlepage}
\begin{flushleft}
\vglue 1.5 true cm

\hspace{3mm}\textsc{\Huge Freie}

\hspace{3mm}\textsc{\Huge Software} 
\newcommand{\HRule}{\rule{1\linewidth}{0.5mm}}
\HRule \\[0.4cm]
\hspace{3mm}\textsc{\textbf{\Huge #1}}

\vglue 1 true cm

\begingroup
\leftskip3mm
{ \large\it
Acht Freie \\
Computerprogramme, \\
die Deine Rechte \\
als Nutzer/-in \\
respektieren \\
und schützen. \\
}
\endgroup

\end{flushleft}
\vfill
\hglue 2 true cm \includegraphics[scale=0.25]{cwfs-bw}
\vglue 1 true cm
\begin{flushright}
\includegraphics[width=5cm]{logo_muc_plain}
\end{flushright}
\end{titlepage}

\flushleft
}


%%% Backpage %%%
\newcommand{\backpage}{
\newpage
\section{Was ist »Freie Software«?}
\textbf{»Freie Software«} sind Computer-Programme,
die Dir die folgenden wesentlichen
\textbf{Freiheiten} gewähren:

\begin{itemize}
\item Die Freiheit, das Programm auf jede Weise und zu jedem beliebigen Zweck zu \textbf{verwenden}.
\item Die Freiheit, das Programm zu \textbf{verstehen} und den eigenen Bedürfnissen \textbf{anzupassen}.
\item Die Freiheit, das Programm \textbf{weiterzugeben}, um Deinen Mitmenschen zu helfen.
\item Die Freiheit, das Programm zu \textbf{verbessern} und diese Verbesserungen zu veröffentlichen.
\end{itemize}

Dazu wird die Software unter einer \textbf{Freien Software-Lizenz}
– wie der General Public License (GPL) – angeboten
und der \textbf{Quelltext} zur Verfügung gestellt.

Während die einmal erstellte Software damit in aller Regel auch kostenlos ist,
leben die Ersteller/-innen zum Beispiel von Auftragsarbeiten, Support-Verträgen mit Firmen
oder auch Spenden der glücklichen Nutzer/-innen.

Die hier aufgeführten Programme sind nur
eine winzige \textbf{Auswahl} der verfügbaren Freien
Software, um Dir die erste Annäherung so einfach und sicher wie möglich
zu gestalten. Weitere Informationen zu Freier Software,
Freien Software-Lizenzen und dem »Free Software Movement«
findest Du unter \url{https://fsfe.org}.

\vglue  1 true cm
\textcolor{darkgray}{\scriptsize
\textbf{Impressum:} \ 
Christian Kalkhoff  \ | \ Riesstraße 60  \ | \  80993 München \\
}
}
